%! Modeling with Quadratics — Unit 2, Lesson 2.7 — Guided Notes (Answer Key)
\documentclass[10pt]{article}
\usepackage{algebra2-article}
\usepackage{algebra2-key}   % pulls in -boxes; do NOT also load -boxes
\newcommand{\TallMath}[1]{$\displaystyle #1\rule[-1.4em]{0pt}{3.2em}$}
\newcommand{\vocabans}[2]{%
  \par\noindent\textbf{\textcolor{forest}{#1:}}\\[1pt]\ansline{#2}\par}

\begin{document}

\pageheader{Unit 2, Lesson 2.7}{Guided Notes --- Answer Key}

\vspace{0.06in}

\begin{objectivebox}
\textbf{By the end of this lesson, I will be able to\dots}
\begin{itemize}\setlength{\itemsep}{2pt}
  \item build a quadratic model from a projectile, area, or revenue situation;
  \item read the $y$-intercept as the start, a zero as ``reaches $0$,'' and the \emph{vertex} as the max/min;
  \item interpret every answer in context with units, and reject values that can't happen.
\end{itemize}
\end{objectivebox}

\vspace{0.06in}

\begin{vocabbox}
\small
Fill in each term as we build it together.
\par\vspace{2pt}
\vocabans{Quadratic model}{a quadratic function used to describe a real situation (height, area, revenue)}
\vocabans{Projectile model $h(t)=-16t^2+v_0t+h_0$}{$h_0=$ launch height, $v_0=$ launch speed; feet and seconds}
\vocabans{Maximum / minimum value (at the vertex)}{the output at the vertex --- the greatest/least the model reaches}
\vocabans{Optimize}{find the input giving the largest (or smallest) value --- the vertex}
\vocabans{Feasible (reasonable) domain}{the inputs that make sense in context (e.g.\ $t\ge0$, a width $>0$)}
\end{vocabbox}

\begin{hookbox}
You can now graph, factor, and solve any quadratic. The last step is the most useful: \emph{building} one from a real
story --- a ball in the air, a pen you're fencing, a price you're setting --- and then letting the graph's
\textbf{features} answer the question. The whole trick is a short map: \textbf{start} $\to$ $y$-intercept,
\textbf{``reaches $0$''} $\to$ a zero, \textbf{``the most / the best''} $\to$ the vertex.
\end{hookbox}

\boxguard[30]
\begin{notesbox}{1. Projectile motion --- one model, three questions}
An object thrown or dropped near the ground follows $\ h(t)=-16t^2+v_0t+h_0\ $ (height in feet, $t$ in seconds),
where $h_0$ is the \textbf{launch height} and $v_0$ the \textbf{launch speed}. A ball is thrown \emph{up} at $32$ ft/s
from a $48$-ft ledge: $\ h(t)=-16t^2+32t+48$.
\begin{center}
\begin{tikzpicture}[xscale=1.45, yscale=0.052]
  \draw[->, semithick, charcoal] (-0.15,0)--(3.7,0) node[below right, font=\scriptsize]{$t$ (s)};
  \draw[->, semithick, charcoal] (0,-3)--(0,74) node[left, font=\scriptsize]{$h$ (ft)};
  \draw[very thick, forest, smooth, samples=60, domain=0:3] plot (\x,{-16*\x*\x+32*\x+48});
  \draw[dashed, linegray] (1,0)--(1,64)--(0,64);
  \node[circle, fill=charcoal, inner sep=1.4pt] at (0,48){};
  \node[circle, fill=charcoal, inner sep=1.4pt] at (1,64){};
  \node[circle, fill=charcoal, inner sep=1.4pt] at (3,0){};
  \node[font=\scriptsize, anchor=west]  at (1.08,64){max $(1,64)$};
  \node[font=\scriptsize, anchor=west]  at (0.08,50){start $(0,48)$};
  \node[font=\scriptsize, anchor=south west] at (2.75,1){lands $t=3$};
\end{tikzpicture}
\end{center}
\vspace{-2pt}
Match each question to a \textbf{feature}, then find it:
\begin{center}
\renewcommand{\arraystretch}{1.4}
\begin{tabularx}{\linewidth}{@{}p{0.36\linewidth} l X@{}}
\textbf{Question} & \textbf{Feature} & \textbf{Work \& answer} \\
\hline
How high does it start?      & $y$-intercept & $h(0)=\ans{$48$}$ \ ft \\
Highest point, and when?     & vertex        & $t=-\dfrac{b}{2a}=\ans{$1$}$;\ \ $h(t)=\ans{$64$}$ ft \\
When does it hit the ground? & zero ($h=0$)  & $-16t^2+32t+48=0\Rightarrow t=\ans{$3$ (reject $-1$)}$ \\
\end{tabularx}
\end{center}
\emph{Reject} the negative root --- a time before launch can't happen. Max height \ans{$64$} ft at $t=\ans{$1$}$ s; lands at $t=\ans{$3$}$ s.
\end{notesbox}

\begin{notesbox}{2. Maximum area --- the vertex is the ``most''}
A farmer fences a rectangular pen against a straight barn wall, so only \emph{three} sides need fence, and she has
$40$ ft of it. Let the width be $x$.
\begin{center}
\begin{tikzpicture}[scale=0.85]
  \draw[very thick, charcoal] (-0.3,1.6)--(5.3,1.6);
  \node[font=\scriptsize, charcoal, anchor=south] at (2.5,1.65){barn wall (no fence)};
  \draw[very thick, forest] (0,0)--(0,1.6);
  \draw[very thick, forest] (5,0)--(5,1.6);
  \draw[very thick, forest] (0,0)--(5,0);
  \node[font=\scriptsize, forest, anchor=east] at (-0.05,0.8){$x$};
  \node[font=\scriptsize, forest, anchor=west] at (5.05,0.8){$x$};
  \node[font=\scriptsize, forest, anchor=north] at (2.5,-0.05){$40-2x$};
\end{tikzpicture}
\end{center}
The two widths and one length use all the fence: $\ x+x+(\text{length})=40$, so length $=\ans{$40-2x$}$. \\[3pt]
Area: $\ A(x)=x(40-2x)=\ans{$-2$}\,x^2+\ans{$40$}\,x$. \\[3pt]
The pen is largest at the \textbf{vertex}: $\ x=-\dfrac{b}{2a}=\ans{$10$}$, so $A=\ans{$200$}$ ft$^2$, with
dimensions $\ans{$10$}\times\ans{$20$}$ ft. \\[3pt]
\textbf{Feasible domain:} $x$ must satisfy $\ans{$0<x<20$}$ (a width can't be $0$ or use more than the fence allows).
\end{notesbox}

\begin{notesbox}{3. Optimize revenue --- best price}
A theater sells $200$ tickets at \$8 each. A survey says each \$1 \emph{increase} in price will lose $20$ buyers.
Let $x=$ the number of \$1 increases.
\[
  \text{price}=(8+x)\text{ dollars},\qquad \text{attendance}=(200-\ans{$20$}\,x)\text{ people}.
\]
Revenue $=$ price $\times$ attendance:
\[
  R(x)=(8+x)(200-20x)=\ans{$-20$}\,x^2+\ans{$40$}\,x+\ans{$1600$}.
\]
Best price is an \textbf{optimization} --- the vertex: $\ x=-\dfrac{b}{2a}=\ans{$1$}$. So charge
\$$\ans{$9$}$, sell $\ans{$180$}$ tickets, for a maximum revenue of \$$\ans{$1620$}$.
\end{notesbox}

\begin{notesbox}{The strategy (use it every time)}
\small
\textbf{(1)} Define the variable (units). \quad \textbf{(2)} Write the quadratic model. \quad \textbf{(3)} Decide the
feature the question wants --- \emph{start} $\to$ $y$-int, \emph{``reaches $0$''} $\to$ zero, \emph{``most/least/best''}
$\to$ vertex. \quad \textbf{(4)} Find it (Lessons 2.1--2.5). \quad \textbf{(5)} Interpret with units; reject
impossible values.
\end{notesbox}

\begin{practicebox}
\textbf{Build the model, then answer the question. Include units.}
\begin{enumerate}[leftmargin=1.9em, itemsep=4pt]
  \item A firework launches from the ground: $h(t)=-16t^2+64t+80$. Find its \emph{maximum height} and \emph{when it lands}.
    \begin{work}
      t &= -\frac{64}{2(-16)} = 2\text{ s} \\
      h(2) &= -16(4)+64(2)+80 = 144\text{ ft} \\
      -16t^2+64t+80 &= 0 \\
      -16(t-5)(t+1) &= 0 \\
      t &= 5\text{ s} \quad\text{(reject } t=-1)
    \end{work}
  \item A dog run against a wall uses $60$ ft of fence on three sides. What width $x$ gives the \emph{maximum area}, and what is that area?
    \begin{work}
      A(x) &= x(60-2x) = -2x^2+60x \\
      x &= -\frac{60}{2(-2)} = 15\text{ ft} \\
      A(15) &= 15(30) = 450\text{ ft}^2
    \end{work}
  \item A shop's daily revenue is $R(x)=-10x^2+300x$ ($x$ in \$1 price cuts). What price change \emph{maximizes} revenue, and what is the maximum?
    \begin{work}
      x &= -\frac{300}{2(-10)} = 15 \\
      R(15) &= -10(225)+300(15) = \$2250
    \end{work}
\end{enumerate}
\end{practicebox}

\end{document}
